\magnification=\magstep1
\input notesmac.tex
\nopagenumbers


\centerline{\titlefont CZ3106\ \  Symbolic Computing, Tutorial 4}
\bigskip
\bigskip
\line{A/P~Wang Jian-Sheng \hfill week 10-15 March 03}
\bigskip
\bigskip

%----------------------------------------------------------------------- 
\problem  Give some examples of group, ring, and field.  
Decide the algebraic structures of the followings: 
(a) a set of all integers with addition,  (b) a set of all integers with
addition and multiplication, (c) set of all rational numbers with
addition and multiplication, (d) the set of all real numbers 
with addition and multiplication, (e) the set of integers from
0 to $p-1$ with addition and multiplication modulo $p$,  for any
integer $p$, and for $p$ being a prime. (f)  a set of all
two by two  unitary matrices with multiplication. 


\problem  Quaternions are the elements of an associative algebra of rank
4 with the basis elements
$$\bf 1,\ i,\ j,\ k$$
where ${\bf 1}$ is the identity elements, i.e., 
$\bf 1 i = i$, $\bf 1 j = j$, $\bf 1 k = k$.
The compositions (multiplications) of two basis elements are
$$ {\bf i}^2 = {\bf j}^2 = {\bf k}^2 = - {\bf 1}$$ 
and 
$$\bf i j = k,\ jk = i,\ ki = j,\ ji = -k,\ kj = -i,\ ik = - j.$$
A quaternion is represented as (where $a$, $b$, $c$, and $d$ are real numbers)
$$ a {\bf 1} + b\,{\bf i} + c\,{\bf j} + d\,{\bf k}.$$
Two quaternions can be added, where the corresponding coefficients are 
simply added.  Two quaternions can be multiplied, with the usual distribution
law of multiplication over addition and the rules for the multiplications
of the basis above.  
\item{} Do quaternions form a group, a ring, an integral domain, or
field?  The answer depends on the operation (addition or multiplication)
and can be multiple possibilities.

\problem consider a set of all polynomials of real coefficients
modulo $(1+x^2)$.  That is all $r(x)$ such that
$$ P(x) = Q(x)*(1+x^2) + r(x),\quad  {\rm deg}(r(x)) < 2$$
With addition and multiplication modulo $(1+x^2)$ defined on the
polynomials, show that it forms a field.  Show that it is a linear
vector space of dimension 2.  Show that the field is isomorphic to
a set of complex numbers.  [This question is slightly ahead of
the lecture].

\bye
